\chapter{NUMA、互连与缓存一致性}

多核心 CPU 不是把多个核心简单放在一起。核心之间通过片上互连通信，共享某些缓存层级，连接内存控制器，并用缓存一致性协议维护共享内存语义。多 socket 机器还会引入 NUMA：不同核心访问不同内存节点的成本不同。

\section{从共享内存到一致性流量}

在 C++ 程序里，多个线程似乎可以直接读写同一地址。但硬件上，写一个共享 cache line 需要让其他核心的副本失效或更新。一个原子自增不是普通加法，它需要获得 cache line 独占权，并按内存序约束发布结果。争用越激烈，cache line 在核心之间来回迁移越频繁。

这解释了为什么全局计数器扩展性差。四个线程分别做本地计数，最后合并，通常比四个线程同时 \code{fetch_add} 同一个原子变量快。前者减少共享写，后者让每次更新都进入一致性协议热路径。

\topic{MESI 的直觉}

很多一致性协议都可以用 MESI 直觉理解：Modified 表示本核心有被修改且尚未写回的独占副本，Exclusive 表示本核心有未修改的独占副本，Shared 表示多个核心可以读，Invalid 表示副本无效。当一个核心要写 Shared 行时，需要让其他核心副本失效；当其他核心随后要读，又要重新获取数据。

软件看不到这些状态名，但能看到它们的后果。一个只读共享表可以被多个核心同时缓存，扩展性很好；一个频繁更新的共享 head 指针会让 cache line 在核心之间跳来跳去；一个包含多个线程本地字段的结构体如果没有 padding，也会因为 false sharing 触发同样的迁移。

真实处理器可能使用 MESIF、MOESI 或其他变体，但教学时先抓住一个核心事实：读共享容易扩展，写共享需要取得所有权。Modified 或 Owned 这类状态的具体名称不重要，重要的是写入会触发状态迁移，状态迁移会占用互连和缓存层级资源。高性能并发设计的第一反应，不应该是“换一条更快的原子指令”，而应该是“能不能减少共享写”。

\begin{lstlisting}[numbers=none]
read mostly table:
  core0 Shared
  core1 Shared
  core2 Shared

hot atomic counter:
  core0 Modified -> core1 Modified -> core2 Modified -> core0 Modified
\end{lstlisting}

上面的文本图表达两种完全不同的共享。前者多个核心一起读，硬件可以复制；后者多个核心轮流写，所有权不断迁移。很多扩展性问题就藏在这条差异里。

\topic{一致性和内存模型的分工}

缓存一致性解决的是“同一内存位置的多个副本如何保持一致”，内存模型解决的是“不同内存操作之间能被哪些线程以什么顺序观察”。硬件一致性不等于程序自动无数据竞争。没有同步的普通读写仍然可能在 C++ 层面形成未定义行为。

一致性协议通常围绕状态机工作，例如 Modified、Exclusive、Shared、Invalid 这类状态。软件不需要手写这些状态，但必须理解状态迁移的成本：共享读便宜，共享写贵，频繁写同一行最贵。

C++ 内存模型和硬件一致性的边界也要分清。硬件一致性通常保证同一 cache line 或同一地址的副本最终按协议协调，但 C++ 程序如果对普通变量并发读写而没有同步，语言层面就是数据竞争。语言规则先决定程序是否有定义，硬件规则再决定有定义的同步怎样落地。不要用“我的 CPU 有缓存一致性”来为未同步普通变量辩护。

\topic{Store buffer 和可见性}

核心执行 store 后，写入可能先进入 store buffer。对本线程来说，后续 load 可能通过 store forwarding 看到自己的写；对其他核心来说，这个写何时可见还受缓存一致性和内存序约束影响。锁释放、release store 和内存屏障会把这种可见性变成可依赖的同步关系。

这就是为什么“我在写 flag 前已经写了 data”在并发中不够。编译器和硬件都可能在允许范围内重排或延迟可见性。正确发布对象需要用 release/acquire、mutex 或其他同步机制建立 happens-before。

Store buffer 还会影响锁的性能。释放锁时，持锁线程在临界区内的写入必须以满足同步语义的方式对后续拿锁线程可见。获取锁时，线程可能需要等待锁所在 cache line 的最新状态。低竞争时这条路径很短；高竞争时，锁变量所在 cache line 会成为所有核心争夺的热点。锁的成本不是一个固定常数，而是竞争、拓扑和写入路径共同决定的结果。

\section{从页面位置到拓扑感知}

NUMA 机器上，内存离某些核心更近。线程在节点 A 上运行，却频繁访问节点 B 分配的内存，会产生远端访问。Linux 的 first-touch 策略让首次触碰页面的线程影响页面放置，因此并行初始化和线程绑定会影响后续计算性能。

NUMA 优化的基本原则是让计算靠近数据。大数组按块分给线程时，初始化也应按相同块并行完成。跨节点共享队列和全局分配器要谨慎，因为它们可能把远端访问和同步争用叠加在一起。

NUMA 不只出现在多 socket 服务器上。有些桌面和移动平台也有非均匀的缓存、内存或芯片let 拓扑，只是操作系统暴露方式不同。教材里使用 NUMA 这个词，是为了让读者形成拓扑意识：核心不是一个无差别集合，内存也不是一个无差别池。线程、数据和 I/O 设备之间的位置关系会影响成本。

\topic{First-touch 和并行初始化}

Linux 常见 first-touch 策略意味着页面通常分配到第一次写入它的线程所在 NUMA 节点。如果主线程单线程初始化一个大数组，然后工作线程分布在多个节点上处理这个数组，很多工作线程可能访问远端内存。更好的方式是让每个工作线程初始化自己后续要处理的分块。

这条规则在 benchmark 中尤其重要。若初始化方式和计算方式不一致，测到的可能是内存放置错误，而不是算法本身。实验报告必须记录初始化策略、线程绑定和 NUMA 策略。

\topic{互连不是无限带宽}

核心、缓存、内存控制器和 I/O 设备之间的互连带宽有限。一个程序即使单线程局部性很好，多线程扩展后也可能撞上内存带宽或互连带宽上限。此时继续增加线程，只会让每个线程分到更少带宽，吞吐不再提高。

互连瓶颈常常表现为“每个线程单独跑都快，一起跑就都慢”。原因可能是多个核心争用同一内存控制器，多个 socket 之间远端访问过多，或者共享写导致 cache line 在互连上来回移动。排查时要把线程数、绑定位置、数据放置和访问模式作为实验矩阵，而不是只改算法代码。

\section{从拓扑记录到分层资源管理}

拓扑感知不是过早优化，而是避免把硬件当作平面资源池。线程池可以按 NUMA 节点分组，内存池可以按节点分配，队列可以按生产者或消费者分片，归约可以先在 socket 内合并再跨 socket 合并。这样的层次化设计和分布式系统中的本地聚合很像：先减少近处流量，再把少量结果发到远处。

\topic{Linux 观察入口}

在 Linux 上，拓扑可以从 \cmd{lscpu}、\filepath{/sys/devices/system/cpu/}、\filepath{/sys/devices/system/node/} 和 \cmd{numactl --hardware} 开始观察。线程绑定可以用 \cmd{taskset}、\cmd{numactl --physcpubind} 或程序内亲和性 API。内存策略可以用 \cmd{numactl --membind}、\cmd{numactl --interleave} 和 first-touch 实验观察。若权限允许，\cmd{perf c2c} 可以帮助识别 cache-to-cache 流量，\cmd{perf stat} 的 cache、LLC、context-switch 和迁核指标也能提供间接证据。

源码阅读可以从调度拓扑、NUMA 策略和内存分配路径切入。不要试图一次看完整调度器或完整内存管理。更好的问题是：线程迁移时哪些状态会变化，first-touch 影响页面放置的路径在哪里，NUMA 策略怎样被记录和查询，某个 benchmark 的现象能不能从这些路径得到合理解释。

\topic{层次化归约案例}

并行归约是理解拓扑的好例子。最差设计是所有线程对同一个原子变量累加。改进设计是每个线程写自己的 partial，最后一个线程合并。进一步的拓扑感知设计是每个核心或每个 worker 先本地累加，每个 NUMA 节点内再合并，最后跨节点合并。这样做减少共享写，也减少远端通信。

这个结构和分布式计算中的树形聚合完全一致。不同只是成本单位变化：单机里远处是另一个 NUMA 节点，分布式里远处是另一台机器。理解这条相似性，可以帮助读者把第二册前后内容串起来。

\topic{一致性流量的三个等级}

共享内存程序里的数据共享可以分成三个等级。第一个等级是只读共享，例如配置表、只读索引、常量权重、不可变输入数组。多个核心可以同时缓存这些数据，扩展性通常很好。第二个等级是低频写共享，例如阶段结束时写一次完成标志、每秒汇总一次指标、偶尔替换一份配置快照。它需要同步，但不一定成为瓶颈。第三个等级是高频写共享，例如所有请求更新同一个原子计数器、所有 worker 争用同一个队列 tail、多个线程反复写同一 cache line 上的不同字段。这个等级最危险，往往表现为核心数越多越慢。

优化时要先判断共享等级，而不是先决定使用 mutex、atomic 还是无锁结构。mutex 保护低频复杂状态可能完全合理；atomic 保护高频热点计数器仍然可能很慢；无锁队列如果所有生产者都争用同一个 tail，也逃不开一致性流量。把共享写变成分片写、批量写、线程本地写或按 NUMA 节点写，通常比替换同步原语更根本。

这也是第二册区分“正确性同步”和“扩展性设计”的原因。正确性同步回答的是数据竞争、可见性和不变量；扩展性设计回答的是共享写频率、cache line 迁移和拓扑距离。一个程序可以同步正确但扩展性很差，也可以局部很快但同步语义错误。高质量系统必须同时满足二者。

\topic{目录协议和窥探协议的直觉}

一致性协议在小规模系统中可以用广播窥探思想理解：一个核心想写某条 cache line，就向其他核心广播或通过互连通知，让其他副本失效。核心数量增加后，纯广播会变贵，因此很多系统使用目录式信息，记录某条 cache line 可能在哪些核心或节点有副本，再有针对性地发送消息。具体实现非常复杂，但软件层面的直觉很简单：共享范围越大，维护一致性的通信越难便宜。

这条直觉能解释为什么“全局共享结构”在小机器上看不出问题，在大机器上突然变差。四核心机器上，一个全局队列的 cache line 迁移还勉强能承受；多 socket 服务器上，迁移可能跨互连和 NUMA 节点，延迟和带宽成本都更高。工程设计不能只在开发机上看正确性测试通过，还要在目标拓扑上看扩展曲线。

目录和窥探的差别不需要读者背协议细节，但要形成层次化思维。让数据只在一个核心私有，成本最低；让数据在一个 socket 内共享，成本较低；让数据跨 socket 高频写共享，成本高；让数据跨机器共享，成本更高。线程本地 partial、socket 内合并、跨 socket 合并、跨机器 reduce，就是把这个成本阶梯显式写进算法。

\topic{内存控制器和带宽饱和}

NUMA 节点通常连接自己的内存控制器。一个节点上的线程大量访问本地内存，会消耗本地控制器带宽；远端访问则需要经过互连到另一个节点的内存控制器。若所有线程都访问一个节点上的内存，即使 CPU 核心分布在多个节点上，带宽也可能集中压在一个控制器上。此时线程数增加，吞吐未必提高，因为瓶颈不是核心数量，而是某个内存节点的带宽。

first-touch 实验的价值就在这里。主线程初始化全部页面，可能让所有页面落在一个 NUMA 节点；worker 分布到多个节点计算时，远端访问集中出现。按 worker 分块初始化，可以让页面分散到各自节点，使多个内存控制器共同工作。若实验显示 parallel first-touch 版本带宽更高，就说明优化来自页面放置和控制器并行，而不是加法本身改变。

带宽饱和也有反直觉现象。使用更多线程可能让总吞吐略增，但每线程吞吐下降；使用 SMT 可能让内存延迟隐藏更好，也可能只增加争用；把线程平均分到两个节点可能提高总带宽，也可能因为跨节点合并阶段变慢。结论必须来自实验矩阵，不应写成固定经验。

\topic{拓扑感知队列}

队列是互连成本的集中点。一个全局 MPMC 队列很容易成为共享写热点：所有生产者争用 tail，所有消费者争用 head 或槽位状态。低线程数下它简单好用，高线程数下它可能把所有 worker 绑在一组 cache line 上。拓扑感知设计通常会拆成多级队列：每个 worker 或每个 NUMA 节点有本地队列，只有负载不均时才访问其他队列。

线程池的工作窃取就是这个思想。worker 优先使用本地 deque，空闲时才偷别人的任务。NUMA 感知线程池还可以先在同节点 worker 中窃取，再跨节点窃取。这样做不是为了复杂而复杂，而是为了让高频路径局部化、低频路径承担远端成本。

有界队列的背压也要考虑拓扑。若所有生产者把任务推到单个队列，队列满时所有生产者都等待同一个条件变量，唤醒和竞争集中；若每个分片有自己的队列，热点分片仍会背压，但无关分片可以继续执行。分片设计会增加任务路由和负载均衡复杂度，因此要先测出全局队列确实是瓶颈，再拆分。

\topic{Linux NUMA 策略和迁移}

Linux 提供多种 NUMA 相关策略：默认 first-touch、按节点绑定、交错分配、自动 NUMA balancing 等。自动机制可以在一般场景改善局部性，但它不是性能实验的替代品。自动迁移页面本身有成本，统计窗口和迁移决策也可能与 benchmark 生命周期不匹配。短时间 benchmark 尤其容易受到初始化策略影响。

用户态实验应记录三类信息。第一，线程位置：线程是否绑定，绑定到哪些 CPU，是否发生迁移。第二，页面位置：是否通过 first-touch、membind 或 interleave 控制，是否能观测各节点内存使用。第三，访问形状：线程是否主要访问本地分块，是否有跨节点共享写，合并阶段是否集中到一个节点。只记录 \cmd{numactl --hardware} 不够，必须把拓扑和程序行为联系起来。

源码阅读可以围绕 \filepath{mm/mempolicy.c}、NUMA balancing 相关路径和调度拓扑展开。不要试图一次读懂所有内存管理。合格阅读报告应选择一个问题，例如“为什么主线程初始化会影响页面节点”，然后追踪 VMA 策略、缺页处理和页面分配之间的关系。

\topic{C++ 接口如何表达拓扑}

C++ 标准库不直接暴露 NUMA 拓扑，但工程代码可以在接口层表达位置意识。例如线程池可以接收 worker 拓扑描述；内存池可以按节点创建 arena；任务可以带有 preferred node；归约可以返回分层 partial；队列可以按 shard id 提交。接口不需要一开始绑定某个 Linux API，但要给后续拓扑策略留下位置。

下面是一个概念性的任务描述。

\begin{lstlisting}[language=C++,caption={任务元数据中保留位置偏好}]
struct TaskPlacement {
    std::int32_t preferred_worker = -1;
    std::int32_t preferred_numa_node = -1;
};

struct ComputeTask {
    std::size_t begin = 0;
    std::size_t end = 0;
    TaskPlacement placement;
};
\end{lstlisting}

这段代码不是要读者立刻实现完整 NUMA runtime，而是提醒接口设计不要把任务看成无位置的函数。对数组块任务，\code{begin/end} 暗含数据位置；对分布式 shuffle，partition id 暗含目标节点；对线程池，本地队列暗含 worker 位置。位置语义一旦在接口中完全丢失，后面很难再补回来。

\topic{拓扑不是核心数列表}

很多程序只读取“有多少 CPU”，然后创建同样数量线程。这不足以描述现代机器。CPU 拓扑包含 socket、NUMA node、core、SMT sibling、cache 共享层级、内存控制器和 I/O 设备位置。两个逻辑 CPU 可能是同一个物理核心的 SMT sibling，也可能在同一个 socket 的不同核心，也可能跨 socket。它们之间共享资源和通信成本完全不同。

性能报告应至少记录 \cmd{lscpu} 中的核心数、线程数、socket 数、NUMA 节点和 cache 信息。在多 socket 机器上，还应记录每个 NUMA node 的 CPU 列表。线程池设计也要区分 worker id 和 CPU id。worker id 是运行时内部编号，CPU id 是操作系统调度对象，NUMA node 是内存位置。把它们混为一谈，会让绑定和 first-touch 实验失去意义。

\topic{互连成本如何出现}

多核心之间通过片上网络、环、mesh 或 socket 间互连通信。软件不会直接调用“互连”，但互连成本会出现在 cache line 迁移、远端内存访问、跨 socket 锁竞争和共享 LLC 争用中。一个全局原子计数器在单 socket 小机器上看起来还能接受，在多 socket 上可能因为 cache line 在 socket 间来回迁移而急剧变慢。

互连成本也会出现在读多写少结构中。只读共享通常可扩展，因为多个核心可以缓存副本；写入会让其他副本失效。若一个配置对象每秒更新一次，问题不大；若一个统计字段每次请求都写，问题很大。软件层面要区分只读共享、低频写共享和高频写共享。这比盲目选择 atomic 或 mutex 更重要。

\topic{目录协议的工程影响}

目录协议维护某条 cache line 的副本位置或所有权信息。软件工程师不需要实现目录协议，但要理解它的结果：共享范围越大，一致性维护越贵；写者越分散，所有权迁移越频繁；数据越能保持本地私有，扩展性越好。分片、线程本地、per-NUMA partial 和分层合并，都是把共享范围缩小。

这也解释为什么“每个线程写自己的数组下标”不一定安全高效。如果这些下标落在同一 cache line，硬件仍然认为它们共享同一个一致性单位。软件的逻辑独立必须映射成硬件的物理独立，才真正减少一致性流量。

\topic{First-touch 的细节}

First-touch 通常发生在页面第一次被写入或实际分配物理页时。只调用 \code{reserve} 或 \code{malloc} 可能只获得虚拟地址，不一定分配物理页。第一次写入触发缺页，内核选择物理页所在 NUMA 节点。若主线程初始化全部数组，页面可能集中在主线程所在节点；后续其他节点线程访问就变成远端访问。

因此 NUMA 实验要把分配、初始化和计算分开。分配阶段只申请地址；初始化阶段按目标线程分块写入，建立页面位置；计算阶段使用同样分块读取。若初始化也计入总时间，要单独报告，因为 parallel first-touch 可能把初始化变快或变慢。若只关心计算阶段，就要在计时前完成初始化并预热。

\begin{lstlisting}[numbers=none]
NUMA experiment phases:
  allocate virtual memory
  initialize pages by owner worker
  optionally warm up
  compute with the same worker-to-range mapping
  merge partials by node then globally
\end{lstlisting}

\topic{线程绑定的收益和风险}

线程绑定可以减少调度迁移，让 worker 更稳定地访问本地 cache 和本地 NUMA 内存。它也可能降低操作系统调度灵活性。如果绑定策略错误，把多个 heavy worker 绑到同一个核心或同一个 NUMA 节点，就会变慢。绑定还要考虑系统上其他进程、容器配额和可用 CPU 集合。

实验中，绑定的价值是减少变量。未绑定实验可能因为操作系统迁移而样本波动；绑定后更容易解释线程和数据位置。但生产系统是否绑定，要看服务类型和部署环境。批处理计算通常更适合显式绑定；通用服务可能需要给调度器更多空间。教材应要求读者记录绑定策略，而不是强制所有程序绑定。

\topic{分层归约}

NUMA 友好的归约通常分层。每个 worker 计算本地 partial；同一 NUMA 节点内先合并成 node partial；最后跨节点合并。这样高频读写发生在本地，跨节点通信只在较小 partial 上发生。这个结构也对应分布式系统中的 local combine、node-level reduce 和 global reduce。

\begin{lstlisting}[numbers=none]
hierarchical reduce:
  worker partials: no sharing during hot loop
  node partials: merge within NUMA node
  global result: merge node partials
\end{lstlisting}

分层合并的代价是代码复杂度和额外阶段。若数据很小或机器只有单 NUMA 节点，收益不明显。实验要比较平坦合并和分层合并，并记录每阶段耗时。不要把 NUMA 策略写成不可关闭的复杂路径。

\topic{跨节点队列的代价}

全局队列在 NUMA 机器上常成为跨节点热点。生产者和消费者分布在不同节点，队列锁、head/tail 和槽位状态会在节点间迁移。分片队列或每节点队列可以减少这种迁移。任务优先进入本节点队列；本节点空闲时先偷同节点任务，再跨节点偷。这个策略和分布式调度中的数据本地性一致。

有界队列的容量也可以按节点分配。若每个 NUMA 节点有自己的输入队列，上游可以根据数据位置提交任务；某个节点过载时，其他节点不必被同一个全局队列阻塞。缺点是负载均衡更复杂，任务路由需要知道位置。设计要看工作负载是否有明显数据本地性。

\topic{远端释放和内存池}

NUMA 不只影响分配，也影响释放。一个对象在节点 A 分配，由节点 B 释放，可能导致分配器元数据和 cache line 在节点间移动。高性能分配器常使用线程本地或节点本地缓存，但跨线程释放仍然需要特殊路径。对象池设计应考虑 owner：对象最好由分配它的线程或节点回收，跨节点释放可以先放入 remote free 队列，再由 owner 批量回收。

无锁数据结构的回收协议也会遇到这个问题。Hazard pointer 或 epoch 延迟释放后，最终由哪个线程 delete 节点？如果所有删除都集中到一个线程，可能造成远端释放和内存热点。内存管理和拓扑不能完全分开。

\topic{I/O 设备和 NUMA}

网卡、NVMe 和 GPU 等设备通常连接到某个 NUMA 节点附近。线程处理设备 I/O 时，如果运行在远端节点，可能增加 PCIe 和内存访问成本。高性能网络和存储系统常把中断、队列、缓冲区和处理线程放在同一 NUMA 节点。虽然本册第二册不深入设备驱动，但要让读者知道 I/O 也有拓扑。

分布式 shuffle 如果在真实服务器上运行，网络接收线程、解压线程和 reduce worker 的位置会影响吞吐。若网卡在节点 0，接收缓冲在节点 0，但 reduce worker 在节点 1 读取，远端内存流量会增加。拓扑意识不只属于计算内核，也属于 I/O 管线。

\topic{自动 NUMA balancing}

Linux 自动 NUMA balancing 会采样页面访问，并可能迁移页面到更常访问它的节点。它能改善一些长期运行程序，但也会引入额外 page fault、迁移成本和实验不确定性。短 benchmark 可能还没等自动机制收敛就结束；长服务可能在负载变化时发生页面迁移。

实验报告应记录是否启用自动 NUMA balancing，至少要知道它可能影响结果。若要做可控实验，可以使用 numactl、线程绑定和明确 first-touch。若要评估生产行为，则应在接近真实配置下测量，并记录自动迁移指标。不要让内核自动策略成为看不见的变量。

\topic{容器和 cgroup 的影响}

现代部署常在容器里运行。容器可能限制 CPU 集合、内存节点、CPU 配额和内存上限。程序看到的 \cmd{std::thread::hardware_concurrency} 未必等于实际可用核心数；cgroup 配额可能让线程周期性被 throttled；cpuset 可能只允许访问部分 CPU。性能报告若忽略容器限制，就可能误判扩展曲线。

在 Termux、手机或云环境中，也要承认设备限制。手机通常没有多 socket NUMA，但有大小核、温度和频率限制；云主机可能共享物理资源；容器可能禁止读取某些性能事件。教材项目应记录环境，而不是假设每个读者都有理想服务器。

\topic{拓扑感知不等于硬编码}

接口应表达拓扑信息，但不应把某台机器的拓扑硬编码进算法。正确做法是运行时发现或配置拓扑，然后把任务、内存池和队列映射到抽象 node。算法只知道 preferred node 或 shard id，不直接写死 CPU 编号。这样同一代码可以在单 socket、双 socket、容器和手机上运行，只是策略不同。

Compute Systems Engine 可以从简单配置开始：\code{worker_count}、是否绑定、每个 worker 的 node id。没有 NUMA 信息时，所有 worker 属于 node 0。后续在支持的 Linux 机器上再读取真实拓扑。这样项目不会因为缺少硬件而无法运行，也不会放弃拓扑扩展。

\section{实验边界、降级和项目落地}

没有 NUMA 机器时，不能声称验证了远端内存优化；但仍可以验证 false sharing、线程绑定、任务粒度和内存带宽曲线。单 socket 上 parallel first-touch 和主线程初始化可能差不多，这是正常结果，不是实验失败。报告要写清：本环境只能验证某些机制，不能验证跨 socket 页面放置。

有 NUMA 机器时，也不能只跑一次。需要比较绑定与未绑定、主线程初始化与 worker 初始化、平坦合并与分层合并、单节点内存与 interleave 内存。只有多组对照一致，才能说 NUMA 放置参与了性能结果。

\topic{案例：双 socket 归约}

在双 socket 机器上，大数组归约可以展示 NUMA 的完整链条。版本一由主线程初始化数组，然后所有 worker 分块求和。版本二由每个 worker 初始化自己的块，再求和。版本三按 socket 分组：每个 socket 的 worker 求本地 partial，socket 内合并，再跨 socket 合并。若版本二和三明显更好，说明页面放置和分层合并参与了性能。

实验要记录线程绑定。若 worker 在计算阶段被迁移到其他 socket，first-touch 的页面位置和线程位置又不匹配。还要记录数组大小，太小可能全部被 cache 影响，太大才体现内存控制器。这个案例把第 3 章的 page fault、第 4 章的 NUMA、第 13 章的分层 reduce 连在一起。

\topic{案例：跨 socket 队列}

另一个案例是全局任务队列。所有 worker 从同一个队列取任务，队列锁和 head/tail 位于某个内存节点。跨 socket worker 频繁访问，会让队列状态在 socket 间迁移。本地队列加工作窃取可以减少高频跨 socket 访问：worker 优先本地 pop，空闲时才偷取。

实验可以构造大量短任务，让队列成本显著；再构造少量长任务，让队列成本被计算掩盖。这样读者能看到：同一个队列设计在不同任务粒度下表现不同。NUMA 优化不能脱离任务粒度。

\topic{分层资源管理}

NUMA 机器适合分层资源管理。每个节点有本地 worker、本地内存池、本地队列、本地 partial。跨节点只交换压缩后的 partial 或被窃取的任务。这个结构和分布式系统的节点本地聚合非常像。若在单机上已经学会分层，扩展到集群时会更自然。

分层也会增加复杂度。每个层级都需要指标：本地队列长度、跨节点 steal 次数、节点 partial 大小、远端访问比例。没有指标，分层策略可能只是让代码复杂。项目可以先用单层实现，再在实验分支中加入分层，对比收益。

\topic{大小核和 NUMA 的区别}

手机常见大小核，服务器常见 NUMA。二者都意味着执行资源不均匀，但原因不同。大小核是核心性能和能耗不同，同一内存访问拓扑未必像多 socket 那样分节点；NUMA 是内存位置和访问延迟不同，核心本身可能同构。调度策略也不同：大小核更关注把重计算放大核、后台任务放小核；NUMA 更关注线程和内存同节点。

在 Termux 环境中，读者更可能遇到大小核而不是多 socket NUMA。实验报告应避免把大小核现象误称为 NUMA。若线程数增加后性能下降，可能是小核加入、降频、温度或调度迁移，不一定是远端内存。记录环境是第一步。

\topic{拓扑抽象接口}

项目可以定义一个抽象拓扑接口，而不是直接依赖某个 Linux API。接口返回 \code{worker_count}、\code{node_count}、\code{worker_to_node}、是否支持绑定。单 socket 或手机环境返回一个 node；双 socket 服务器返回多个 node。算法根据 \code{node_count} 决定是否启用分层 partial。这样代码能在不同设备上运行。

\begin{lstlisting}[language=C++,caption={拓扑描述的接口形态}]
struct WorkerTopology {
    std::int32_t worker_count = 1;
    std::int32_t node_count = 1;
    std::vector<std::int32_t> worker_to_node;
};
\end{lstlisting}

这个结构不直接包含 CPU id，是为了让教学项目先表达策略，再逐步接入真实 affinity。若绑定不可用，仍然可以使用 \code{worker_to_node} 做逻辑分层实验。接口应允许能力缺失，而不是让项目只能在服务器上运行。

\topic{拓扑指标}

拓扑感知优化需要指标验证。每个 node 的任务数、处理字节、partial 大小、队列长度、跨节点 steal 次数、远端调度次数都应记录。若启用分层后跨节点 steal 很多，说明负载不均或任务划分不好；若某个 node 长期空闲，说明调度没有利用资源；若 node partial 合并很慢，说明最终合并成为瓶颈。

没有这些指标，拓扑策略容易变成不可解释的复杂代码。项目应先输出简单文本摘要，后续再画图。

\topic{拓扑和故障域}

拓扑不只影响性能，也影响故障域。单机里，一个 NUMA 节点过载可能拖慢本地 worker；集群里，一个 rack、zone 或机房故障会影响一组节点。分布式副本放置通常要跨故障域，避免一处故障丢失多数副本。单机拓扑意识可以自然扩展到集群故障域意识：不要把所有关键副本、任务或队列集中在同一个风险位置。

Compute Systems Engine 的本机模拟可以先用 node id 表示逻辑故障域。后续多进程或多机版本中，node id 可以映射到主机、rack 或 zone。调度器在追求本地性的同时，也要避免把所有副本放同一域。性能和可靠性有时会冲突，系统设计要显式权衡。

\topic{拓扑章节总结}

本章的核心不是让读者背 NUMA API，而是建立位置意识。线程有位置，内存有位置，队列有位置，I/O 设备有位置，副本也有位置。位置决定通信成本和故障相关性。没有位置意识的程序在小机器上可能正常，在大机器上会遇到扩展性和可靠性问题。

\topic{拓扑决策表}

做拓扑相关设计时，先问机器是否真的有多个内存节点，是否能设置亲和性，数据是否大到超过 cache，任务是否访问固定数据块，是否有跨节点共享写，是否需要跨故障域放置副本。若没有 NUMA 或数据很小，复杂拓扑策略可能无收益；若数据大且任务固定访问某块，first-touch 和绑定值得实验；若共享写频繁，先分片再谈绑定；若是分布式副本，优先考虑故障域而不是只考虑延迟。

\topic{拓扑策略的降级设计}

拓扑感知代码必须能降级。很多读者会在手机、容器、云虚拟机或普通单 socket 机器上运行实验，程序不一定能读取完整 NUMA 信息，也不一定有权限设置 affinity。如果代码把真实 NUMA 当成必需条件，就会让教材项目难以运行。更好的策略是把拓扑能力分层：能读取真实 node 时使用真实 node；不能读取时退化成一个逻辑 node；能绑定线程时执行绑定；绑定失败时记录失败并继续运行；能控制内存策略时做 first-touch 实验；不能控制时只做普通分块。

降级不等于忽略拓扑。即使只有一个逻辑 node，代码仍然可以保留 worker 到 node 的映射、分层 partial 的接口和队列分片的抽象。这样程序在小机器上正确运行，在大机器上能启用更多策略。教材项目尤其适合这种设计：先让抽象存在，再逐步接入真实硬件能力。读者不会因为没有双 socket 服务器就学不了拓扑意识。

拓扑策略还要写入输出。若报告只写“启用 NUMA 优化”，但实际上 affinity 设置失败，读者会得到错误结论。程序应输出能力检测结果、降级原因和最终策略。例如 \code{node_count} 是 1，\code{binding_enabled} 是 false，\code{binding_error} 是 permission denied，\code{memory_policy} 是 default。这样实验结果才诚实。高性能系统的第一步不是让所有优化都打开，而是知道哪些优化真的生效。

降级设计还有一个好处：它迫使接口表达意图，而不是表达某个平台的偶然细节。算法真正需要的是“这个任务最好靠近哪一组数据”和“这个 worker 属于哪个执行域”，不一定需要知道具体 CPU 编号。平台能力强时，执行域可以映射到真实 NUMA 节点；平台能力弱时，执行域只是逻辑分组。把意图和机制分开，代码才不会因为换设备而到处改。

这种抽象也能保护学习路径。读者先在普通设备上理解任务、数据和执行域的关系，再到服务器上验证真实远端访问成本。先学清楚问题，再打开平台能力，实验结论会稳得多。


\topic{高密度主线补充：从失败推导机制}

NUMA 和一致性章要从一个扩展性曲线突然变平的问题讲起。这一轮补充专门解决两个问题：第一，避免把本章写成概念枚举；第二，让每个概念都能从多线程计数和大数组归约的失败中自然推出。本章的核心失败不是“代码不会写”，而是线程数量增加后，瓶颈从计算转移到 cache line 所有权、互连带宽和远端内存。所以正文的顺序必须始终保持为：先给出具体任务，再写朴素方案，再观察失败信号，再引入机制，最后给出实验和边界。
这样的写法比直接给定义慢一些，但它更接近工程学习的真实路径。真实系统很少把问题包装成选择题；它只会表现为吞吐上不去、CPU 利用率怪、结果偶尔错、队列越积越多、重启后状态对不上、线上延迟突然变长。读者要学会从这些现象反推机制，而不是看到术语才想起术语。

\topic{一致性所有权}

先把问题收窄到一个可推导的场景：多个核心写同一个计数器时，为什么每次加法都不是普通加法。在多线程计数和大数组归约里，这个问题不会以术语形式出现，而是以结果不稳定、吞吐不增长、延迟抖动或恢复困难的形式出现。如果一上来只背一致性所有权的定义，读者很容易把它当成孤立知识点；更好的顺序是先看朴素方案为什么合理，再看它在真实约束下怎样失败。

朴素方案通常是：所有线程对同一个 atomic 做 fetch add。这个方案的吸引力在于它贴近单线程直觉，代码短，状态少，测试也容易写。但是它隐含了几个没有说出口的假设：输入总是干净，资源近似无限，执行不会交错，失败不会发生，硬件成本和源码结构大体一致。一旦这些假设被打破，cache line 在核心间迁移，吞吐随线程数上升很快停止增长。这时继续在原方案上加 if 或加锁，往往只会把真正的问题埋得更深。

一致性所有权就是从这个失败里长出来的概念。它的核心模型可以这样理解：写共享需要取得 cache line 独占所有权，读共享和写共享成本完全不同。这个模型必须同时放在三个层次看。源码层要说明谁读谁写、谁拥有对象、哪个状态允许变化；运行时层要说明线程、队列、任务和 I/O 怎样排队；硬件或系统层要说明缓存、页、调度、文件系统或网络怎样参与成本。三层缺一层，解释都会变形：只有源码层会看不见隐藏成本，只有硬件层会丢失业务语义，只有运行时层又容易把正确性和性能混在一起。

观察方法不能停在“跑一下变快了”。这一点在本章尤其重要：比较全局 atomic、分片计数器和节点内合并。实验要有 reference，要固定输入规模，要记录环境，要把冷启动、预热、核心循环、提交和清理分开。如果工具权限不足，也要写清楚不能观察什么，而不是把猜测写成结论。系统教材的严谨性来自证据链：现象、假设、对照、指标、结论边界。

一致性所有权也有边界。atomic 保证语言语义，不保证扩展性。工程判断不是把一个技术推到所有地方，而是知道它在哪些约束下成立。读者在本章应形成一种习惯：每学到一个机制，都要问它解决了什么失败、引入了什么新成本、需要什么测试、在什么输入或部署环境下会失效。

\topic{MESI 直觉}

先把问题收窄到一个可推导的场景：只读表和热点计数器为什么都是共享却表现不同。在多线程计数和大数组归约里，这个问题不会以术语形式出现，而是以结果不稳定、吞吐不增长、延迟抖动或恢复困难的形式出现。如果一上来只背MESI 直觉的定义，读者很容易把它当成孤立知识点；更好的顺序是先看朴素方案为什么合理，再看它在真实约束下怎样失败。

朴素方案通常是：把共享数据一概视为危险。这个方案的吸引力在于它贴近单线程直觉，代码短，状态少，测试也容易写。但是它隐含了几个没有说出口的假设：输入总是干净，资源近似无限，执行不会交错，失败不会发生，硬件成本和源码结构大体一致。一旦这些假设被打破，读共享可复制，高频写共享要反复失效或迁移。这时继续在原方案上加 if 或加锁，往往只会把真正的问题埋得更深。

MESI 直觉就是从这个失败里长出来的概念。它的核心模型可以这样理解：Modified、Exclusive、Shared、Invalid 可以帮助理解状态迁移。这个模型必须同时放在三个层次看。源码层要说明谁读谁写、谁拥有对象、哪个状态允许变化；运行时层要说明线程、队列、任务和 I/O 怎样排队；硬件或系统层要说明缓存、页、调度、文件系统或网络怎样参与成本。三层缺一层，解释都会变形：只有源码层会看不见隐藏成本，只有硬件层会丢失业务语义，只有运行时层又容易把正确性和性能混在一起。

观察方法不能停在“跑一下变快了”。这一点在本章尤其重要：用文本状态图解释读多写少和写热点。实验要有 reference，要固定输入规模，要记录环境，要把冷启动、预热、核心循环、提交和清理分开。如果工具权限不足，也要写清楚不能观察什么，而不是把猜测写成结论。系统教材的严谨性来自证据链：现象、假设、对照、指标、结论边界。

MESI 直觉也有边界。真实协议更复杂，教学模型只用于解释软件现象。工程判断不是把一个技术推到所有地方，而是知道它在哪些约束下成立。读者在本章应形成一种习惯：每学到一个机制，都要问它解决了什么失败、引入了什么新成本、需要什么测试、在什么输入或部署环境下会失效。

\topic{Store buffer}

先把问题收窄到一个可推导的场景：写 data 再写 flag 为什么仍需要同步语义。在多线程计数和大数组归约里，这个问题不会以术语形式出现，而是以结果不稳定、吞吐不增长、延迟抖动或恢复困难的形式出现。如果一上来只背Store buffer的定义，读者很容易把它当成孤立知识点；更好的顺序是先看朴素方案为什么合理，再看它在真实约束下怎样失败。

朴素方案通常是：相信源码顺序就是其他线程看到的顺序。这个方案的吸引力在于它贴近单线程直觉，代码短，状态少，测试也容易写。但是它隐含了几个没有说出口的假设：输入总是干净，资源近似无限，执行不会交错，失败不会发生，硬件成本和源码结构大体一致。一旦这些假设被打破，写入可能先在缓冲中，对其他核心的可见顺序必须由同步建立。这时继续在原方案上加 if 或加锁，往往只会把真正的问题埋得更深。

Store buffer就是从这个失败里长出来的概念。它的核心模型可以这样理解：store buffer 和内存模型共同决定发布可见性。这个模型必须同时放在三个层次看。源码层要说明谁读谁写、谁拥有对象、哪个状态允许变化；运行时层要说明线程、队列、任务和 I/O 怎样排队；硬件或系统层要说明缓存、页、调度、文件系统或网络怎样参与成本。三层缺一层，解释都会变形：只有源码层会看不见隐藏成本，只有硬件层会丢失业务语义，只有运行时层又容易把正确性和性能混在一起。

观察方法不能停在“跑一下变快了”。这一点在本章尤其重要：用 release acquire 和错误普通变量版本做对照。实验要有 reference，要固定输入规模，要记录环境，要把冷启动、预热、核心循环、提交和清理分开。如果工具权限不足，也要写清楚不能观察什么，而不是把猜测写成结论。系统教材的严谨性来自证据链：现象、假设、对照、指标、结论边界。

Store buffer也有边界。不要用硬件一致性为 C++ 数据竞争辩护。工程判断不是把一个技术推到所有地方，而是知道它在哪些约束下成立。读者在本章应形成一种习惯：每学到一个机制，都要问它解决了什么失败、引入了什么新成本、需要什么测试、在什么输入或部署环境下会失效。

\topic{NUMA first touch}

先把问题收窄到一个可推导的场景：为什么同一数组在不同初始化方式下多线程带宽不同。在多线程计数和大数组归约里，这个问题不会以术语形式出现，而是以结果不稳定、吞吐不增长、延迟抖动或恢复困难的形式出现。如果一上来只背NUMA first touch的定义，读者很容易把它当成孤立知识点；更好的顺序是先看朴素方案为什么合理，再看它在真实约束下怎样失败。

朴素方案通常是：主线程分配并初始化所有页面。这个方案的吸引力在于它贴近单线程直觉，代码短，状态少，测试也容易写。但是它隐含了几个没有说出口的假设：输入总是干净，资源近似无限，执行不会交错，失败不会发生，硬件成本和源码结构大体一致。一旦这些假设被打破，页面可能集中在一个节点，其他节点 worker 远端访问。这时继续在原方案上加 if 或加锁，往往只会把真正的问题埋得更深。

NUMA first touch就是从这个失败里长出来的概念。它的核心模型可以这样理解：NUMA 把内存位置引入成本模型，first touch 影响页面放置。这个模型必须同时放在三个层次看。源码层要说明谁读谁写、谁拥有对象、哪个状态允许变化；运行时层要说明线程、队列、任务和 I/O 怎样排队；硬件或系统层要说明缓存、页、调度、文件系统或网络怎样参与成本。三层缺一层，解释都会变形：只有源码层会看不见隐藏成本，只有硬件层会丢失业务语义，只有运行时层又容易把正确性和性能混在一起。

观察方法不能停在“跑一下变快了”。这一点在本章尤其重要：比较绑定与未绑定、single touch 与 parallel touch。实验要有 reference，要固定输入规模，要记录环境，要把冷启动、预热、核心循环、提交和清理分开。如果工具权限不足，也要写清楚不能观察什么，而不是把猜测写成结论。系统教材的严谨性来自证据链：现象、假设、对照、指标、结论边界。

NUMA first touch也有边界。没有 NUMA 机器时不能声称验证了远端内存。工程判断不是把一个技术推到所有地方，而是知道它在哪些约束下成立。读者在本章应形成一种习惯：每学到一个机制，都要问它解决了什么失败、引入了什么新成本、需要什么测试、在什么输入或部署环境下会失效。

\topic{互连饱和}

先把问题收窄到一个可推导的场景：每个线程单独跑都快，一起跑为什么都慢。在多线程计数和大数组归约里，这个问题不会以术语形式出现，而是以结果不稳定、吞吐不增长、延迟抖动或恢复困难的形式出现。如果一上来只背互连饱和的定义，读者很容易把它当成孤立知识点；更好的顺序是先看朴素方案为什么合理，再看它在真实约束下怎样失败。

朴素方案通常是：继续增加线程寻找吞吐。这个方案的吸引力在于它贴近单线程直觉，代码短，状态少，测试也容易写。但是它隐含了几个没有说出口的假设：输入总是干净，资源近似无限，执行不会交错，失败不会发生，硬件成本和源码结构大体一致。一旦这些假设被打破，内存控制器、LLC 或 socket 互连达到上限。这时继续在原方案上加 if 或加锁，往往只会把真正的问题埋得更深。

互连饱和就是从这个失败里长出来的概念。它的核心模型可以这样理解：核心、缓存、内存控制器和 socket 之间共享有限互连资源。这个模型必须同时放在三个层次看。源码层要说明谁读谁写、谁拥有对象、哪个状态允许变化；运行时层要说明线程、队列、任务和 I/O 怎样排队；硬件或系统层要说明缓存、页、调度、文件系统或网络怎样参与成本。三层缺一层，解释都会变形：只有源码层会看不见隐藏成本，只有硬件层会丢失业务语义，只有运行时层又容易把正确性和性能混在一起。

观察方法不能停在“跑一下变快了”。这一点在本章尤其重要：画线程数到吞吐曲线，并记录带宽和 LLC 指标。实验要有 reference，要固定输入规模，要记录环境，要把冷启动、预热、核心循环、提交和清理分开。如果工具权限不足，也要写清楚不能观察什么，而不是把猜测写成结论。系统教材的严谨性来自证据链：现象、假设、对照、指标、结论边界。

互连饱和也有边界。带宽瓶颈下更多线程可能只增加争用。工程判断不是把一个技术推到所有地方，而是知道它在哪些约束下成立。读者在本章应形成一种习惯：每学到一个机制，都要问它解决了什么失败、引入了什么新成本、需要什么测试、在什么输入或部署环境下会失效。

\topic{分层归约}

先把问题收窄到一个可推导的场景：为什么先本地合并再跨节点合并比所有线程直接写全局结果更稳。在多线程计数和大数组归约里，这个问题不会以术语形式出现，而是以结果不稳定、吞吐不增长、延迟抖动或恢复困难的形式出现。如果一上来只背分层归约的定义，读者很容易把它当成孤立知识点；更好的顺序是先看朴素方案为什么合理，再看它在真实约束下怎样失败。

朴素方案通常是：最后一个全局变量收集所有 partial。这个方案的吸引力在于它贴近单线程直觉，代码短，状态少，测试也容易写。但是它隐含了几个没有说出口的假设：输入总是干净，资源近似无限，执行不会交错，失败不会发生，硬件成本和源码结构大体一致。一旦这些假设被打破，跨节点高频写或集中合并会制造热点。这时继续在原方案上加 if 或加锁，往往只会把真正的问题埋得更深。

分层归约就是从这个失败里长出来的概念。它的核心模型可以这样理解：分层归约把高频路径留在近处，把远端通信压缩到少量结果。这个模型必须同时放在三个层次看。源码层要说明谁读谁写、谁拥有对象、哪个状态允许变化；运行时层要说明线程、队列、任务和 I/O 怎样排队；硬件或系统层要说明缓存、页、调度、文件系统或网络怎样参与成本。三层缺一层，解释都会变形：只有源码层会看不见隐藏成本，只有硬件层会丢失业务语义，只有运行时层又容易把正确性和性能混在一起。

观察方法不能停在“跑一下变快了”。这一点在本章尤其重要：比较 worker、NUMA node、global 三层耗时。实验要有 reference，要固定输入规模，要记录环境，要把冷启动、预热、核心循环、提交和清理分开。如果工具权限不足，也要写清楚不能观察什么，而不是把猜测写成结论。系统教材的严谨性来自证据链：现象、假设、对照、指标、结论边界。

分层归约也有边界。层级太多会增加阶段开销，小数据未必值得。工程判断不是把一个技术推到所有地方，而是知道它在哪些约束下成立。读者在本章应形成一种习惯：每学到一个机制，都要问它解决了什么失败、引入了什么新成本、需要什么测试、在什么输入或部署环境下会失效。

\topic{拓扑感知队列}

先把问题收窄到一个可推导的场景：全局任务队列为什么在多 socket 上成为热点。在多线程计数和大数组归约里，这个问题不会以术语形式出现，而是以结果不稳定、吞吐不增长、延迟抖动或恢复困难的形式出现。如果一上来只背拓扑感知队列的定义，读者很容易把它当成孤立知识点；更好的顺序是先看朴素方案为什么合理，再看它在真实约束下怎样失败。

朴素方案通常是：所有 worker 共享一个 MPMC 队列。这个方案的吸引力在于它贴近单线程直觉，代码短，状态少，测试也容易写。但是它隐含了几个没有说出口的假设：输入总是干净，资源近似无限，执行不会交错，失败不会发生，硬件成本和源码结构大体一致。一旦这些假设被打破，head、tail、锁和槽位状态在节点间迁移。这时继续在原方案上加 if 或加锁，往往只会把真正的问题埋得更深。

拓扑感知队列就是从这个失败里长出来的概念。它的核心模型可以这样理解：本地队列和工作窃取把高频访问局部化，低频偷取承担远端成本。这个模型必须同时放在三个层次看。源码层要说明谁读谁写、谁拥有对象、哪个状态允许变化；运行时层要说明线程、队列、任务和 I/O 怎样排队；硬件或系统层要说明缓存、页、调度、文件系统或网络怎样参与成本。三层缺一层，解释都会变形：只有源码层会看不见隐藏成本，只有硬件层会丢失业务语义，只有运行时层又容易把正确性和性能混在一起。

观察方法不能停在“跑一下变快了”。这一点在本章尤其重要：记录本地 pop、同节点 steal、跨节点 steal 次数。实验要有 reference，要固定输入规模，要记录环境，要把冷启动、预热、核心循环、提交和清理分开。如果工具权限不足，也要写清楚不能观察什么，而不是把猜测写成结论。系统教材的严谨性来自证据链：现象、假设、对照、指标、结论边界。

拓扑感知队列也有边界。队列分片会增加负载均衡和关闭语义复杂度。工程判断不是把一个技术推到所有地方，而是知道它在哪些约束下成立。读者在本章应形成一种习惯：每学到一个机制，都要问它解决了什么失败、引入了什么新成本、需要什么测试、在什么输入或部署环境下会失效。

\topic{拓扑记录}

先把问题收窄到一个可推导的场景：为什么只写 worker count 不足以复现实验。在多线程计数和大数组归约里，这个问题不会以术语形式出现，而是以结果不稳定、吞吐不增长、延迟抖动或恢复困难的形式出现。如果一上来只背拓扑记录的定义，读者很容易把它当成孤立知识点；更好的顺序是先看朴素方案为什么合理，再看它在真实约束下怎样失败。

朴素方案通常是：报告只记录线程数。这个方案的吸引力在于它贴近单线程直觉，代码短，状态少，测试也容易写。但是它隐含了几个没有说出口的假设：输入总是干净，资源近似无限，执行不会交错，失败不会发生，硬件成本和源码结构大体一致。一旦这些假设被打破，SMT、socket、NUMA node、频率和容器限制都会改变结果。这时继续在原方案上加 if 或加锁，往往只会把真正的问题埋得更深。

拓扑记录就是从这个失败里长出来的概念。它的核心模型可以这样理解：拓扑记录把 CPU id、core、node、cache 共享和绑定策略写进报告。这个模型必须同时放在三个层次看。源码层要说明谁读谁写、谁拥有对象、哪个状态允许变化；运行时层要说明线程、队列、任务和 I/O 怎样排队；硬件或系统层要说明缓存、页、调度、文件系统或网络怎样参与成本。三层缺一层，解释都会变形：只有源码层会看不见隐藏成本，只有硬件层会丢失业务语义，只有运行时层又容易把正确性和性能混在一起。

观察方法不能停在“跑一下变快了”。这一点在本章尤其重要：保存 lscpu、numactl、cpuset 和程序绑定配置。实验要有 reference，要固定输入规模，要记录环境，要把冷启动、预热、核心循环、提交和清理分开。如果工具权限不足，也要写清楚不能观察什么，而不是把猜测写成结论。系统教材的严谨性来自证据链：现象、假设、对照、指标、结论边界。

拓扑记录也有边界。生产系统是否绑定要看部署，不是所有服务都适合强绑定。工程判断不是把一个技术推到所有地方，而是知道它在哪些约束下成立。读者在本章应形成一种习惯：每学到一个机制，都要问它解决了什么失败、引入了什么新成本、需要什么测试、在什么输入或部署环境下会失效。

\topic{案例与实验串联}

案例“全局 atomic 计数器”用于把抽象机制落到一段可复现实验。场景是：多个线程每条记录都更新同一个统计值。这个场景要足够小，小到读者可以手工画出数据流、状态变化和成本路径；又要足够真实，能暴露教材正文讨论的失败模式。

第一版不要急着写最终答案，而应保留一个朴素版本：使用 relaxed fetch add。朴素版本的价值不是性能，而是语义清楚。它告诉我们结果应该是什么，也告诉我们优化前系统在哪些地方偷懒。

第二版再引入改进：改为 per worker 或 per NUMA shard。改进必须说明成本迁移到哪里，而不是只说“更快”。例如减少共享写可能增加最终合并成本，批处理可能增加尾延迟，分片可能引入负载倾斜，异步可能让生命周期更难验证。

验证时重点看：比较扩展曲线和每次操作平均成本。报告里至少写出输入规模、硬件或系统环境、编译选项、运行命令、关键指标和反例。若结果与预期不同，优先检查实验变量，而不是马上改代码。
案例“双 socket 数组归约”用于把抽象机制落到一段可复现实验。场景是：数组大小超过 LLC，worker 分布在两个 socket。这个场景要足够小，小到读者可以手工画出数据流、状态变化和成本路径；又要足够真实，能暴露教材正文讨论的失败模式。

第一版不要急着写最终答案，而应保留一个朴素版本：主线程初始化全部页面。朴素版本的价值不是性能，而是语义清楚。它告诉我们结果应该是什么，也告诉我们优化前系统在哪些地方偷懒。

第二版再引入改进：按 worker 绑定和 first touch 分块初始化。改进必须说明成本迁移到哪里，而不是只说“更快”。例如减少共享写可能增加最终合并成本，批处理可能增加尾延迟，分片可能引入负载倾斜，异步可能让生命周期更难验证。

验证时重点看：观察带宽、远端访问迹象和节点内合并成本。报告里至少写出输入规模、硬件或系统环境、编译选项、运行命令、关键指标和反例。若结果与预期不同，优先检查实验变量，而不是马上改代码。
案例“全局队列与本地队列”用于把抽象机制落到一段可复现实验。场景是：大量短任务由所有 worker 抢占执行。这个场景要足够小，小到读者可以手工画出数据流、状态变化和成本路径；又要足够真实，能暴露教材正文讨论的失败模式。

第一版不要急着写最终答案，而应保留一个朴素版本：一个全局队列存所有任务。朴素版本的价值不是性能，而是语义清楚。它告诉我们结果应该是什么，也告诉我们优化前系统在哪些地方偷懒。

第二版再引入改进：每节点队列加窃取。改进必须说明成本迁移到哪里，而不是只说“更快”。例如减少共享写可能增加最终合并成本，批处理可能增加尾延迟，分片可能引入负载倾斜，异步可能让生命周期更难验证。

验证时重点看：比较短任务和长任务两种粒度下的队列成本。报告里至少写出输入规模、硬件或系统环境、编译选项、运行命令、关键指标和反例。若结果与预期不同，优先检查实验变量，而不是马上改代码。

\topic{Linux 源码阅读与系统观察}

Linux 源码阅读要从用户态现象倒推入口。本章可围绕 \filepath{kernel/sched/topology.c}、\filepath{mm/mempolicy.c}、\filepath{kernel/sched/fair.c}、\filepath{kernel/locking/qspinlock.c} 做窄范围阅读。不要试图一次读完整子系统，先选择一个问题，例如一次阻塞为什么睡眠、一次缺页怎样建立映射、一次唤醒怎样进入 run queue、一次写回怎样变成持久化边界。

阅读报告应固定内核版本、阅读路径和实验命令。第一段写用户态最小程序，第二段写观察到的现象，第三段写源码中的关键结构和状态变化，第四段写结论边界。若当前设备不是对应内核，报告必须说明源码只是机制参考，而不是当前运行系统的直接证据。

源码阅读还要看数据结构，而不只是函数名。队列、树、链表、位图、引用计数、状态枚举、等待队列、页表、inode、bio、task 结构这些细节，决定了机制怎样落地。读者要练习把一个用户态概念翻译成内核对象：线程对应 task，等待对应 wait queue 或 futex 路径，文件缓存对应 page cache，内存策略对应 VMA 和页面分配。

\topic{贯穿项目检查点}

把本章落到贯穿项目时，不要只加一个接口或一个 benchmark。每次新增能力都应有语义目标、最小实现、对照实验、指标和失败注入。项目不是堆功能，而是把本章的机制变成可运行、可检查、可复盘的工程对象。

\begin{lstlisting}[numbers=none]
chapter project checkpoints:
  1. 为计数器实现全局和分片两种路径
  2. 记录 CPU 拓扑和线程绑定
  3. 加入分层归约实验
  4. 统计本地队列和跨节点窃取次数
  5. 写 NUMA 可验证与不可验证边界
\end{lstlisting}

项目报告要回答四个问题。第一，朴素版本是什么，为什么不足。第二，改进版本改变了哪个边界，是数据边界、执行边界、同步边界、I/O 边界还是提交边界。第三，证据是什么，哪些指标支持结论。第四，剩余风险是什么，在哪些输入、机器或失败条件下结论可能不成立。只有这样，项目才不会变成代码展示，而会变成系统能力训练。



\topic{第二轮深水区：把机制讲成可验证能力}

本章继续扩写时，重点不再是补充更多名词，而是把已经出现的机制变成可验证能力。围绕多线程计数和大数组归约，读者最终应能做三件事：解释为什么朴素方案失败，设计能隔离变量的实验，把实验结论落到代码边界。这三件事缺一不可。只会解释会变成空谈，只会跑实验会变成工具使用，只会改代码又会在复杂系统里丢掉语义。
本章的主失败是：线程数量增加后，瓶颈从计算转移到 cache line 所有权、互连带宽和远端内存。这个失败会在不同尺度反复出现。小输入里它可能只是一次结果不稳定；多线程里它可能变成锁竞争、乱序或重复写；I/O 和分布式里它可能变成提交不清、恢复不可靠或重试放大。所以本章的每个机制都要写出尺度变化后的表现。

\topic{深挖：一致性所有权}

围绕“一致性所有权”，先把它放回一个具体问题：多个核心写同一个计数器时，为什么每次加法都不是普通加法。如果这个问题只在本机分布式模拟里出现，读者很容易把它当成局部技巧；但在多线程计数和大数组归约中，它会沿着输入、执行、同步、I/O 和提交一路传播。所以讲解时不能只写定义，而要让读者看到它怎样从一个小失败变成系统性约束。

朴素写法通常是“所有线程对同一个 atomic 做 fetch add”。这句话看似简单，却包含几个默认前提：状态只有一个拥有者，执行顺序可预测，资源足够近，失败不会穿过边界，观察结果能够代表真实成本。这些前提在教学小程序中可能成立，在多线程计数和大数组归约里却会被输入规模、线程交错、硬件拓扑或故障注入逐个打破。

失败信号是“cache line 在核心间迁移，吞吐随线程数上升很快停止增长”。注意，这个失败不一定表现为崩溃。它也可能表现为吞吐不再增长、p99 抖动、重启后结果多了一份、某个分区长尾、CPU 使用率很高但有效工作很少。教材要把这些信号和机制绑定起来，否则读者只会背术语，遇到真实现象时仍然不知道从哪里下手。

更可靠的推导顺序是先写出不变量。对一致性所有权来说，不变量不是一句口号，而是本章要保护的事实：哪些数据只能写一次，哪些状态可以重试，哪些结果可以被缓存，哪些成本必须摊销，哪些错误必须外显。只要不变量没有写清，后面的代码、实验和优化都没有稳定坐标。

然后才引入模型：写共享需要取得 cache line 独占所有权，读共享和写共享成本完全不同。这个模型应同时解释正确性和成本。正确性部分回答“为什么结果可信”，成本部分回答“为什么这个实现会快或慢”。如果一个模型只能解释速度，却解释不了失败恢复，它不适合系统教材；如果只能解释语义，却完全不关心 cache、调度、I/O 或网络成本，也不适合第二册。

实验应围绕“只改变一个变量”设计。针对一致性所有权，最小实验可以保留朴素版本，再实现一个改进版本，输入规模从小到大，线程数或并发度从一到目标机器上限，错误注入从无到有。每一组实验都应记录 reference 结果、核心指标、环境和命令。这样读者看到差异时，可以判断差异来自机制本身，而不是来自初始化、缓存冷热、调度迁移或文件系统状态。

观察手段是：比较全局 atomic、分片计数器和节点内合并。观察不是为了堆工具名，而是为了回答一个具体假设。若假设是共享写导致扩展性下降，就应看线程数曲线、写入布局和 cache 相关指标；若假设是提交边界错误，就应做崩溃点矩阵；若假设是队列背压，就应看水位、等待和上游速率。工具输出必须回到假设，不要把截图或命令输出当成结论。

最后要讲边界：atomic 保证语言语义，不保证扩展性。高质量教材必须把反例放在正文里。读者需要知道什么时候该用这个机制，什么时候它会制造新问题，什么时候应该退回更简单方案。比如一个单线程工具没有并发共享，强行引入复杂运行时只会增加维护成本；一个没有恢复要求的临时分析脚本，也不必承担完整 manifest 协议。

把这一点落实到代码审查，可以要求每个涉及一致性所有权的改动都回答五个问题：朴素版本是什么，新增机制保护了哪个不变量，新增成本在哪里，如何回退，哪些测试证明边界。这五个问题比“用了某某技术”更能判断质量，也能防止团队在没有证据时追逐复杂实现。

\topic{深挖：MESI 直觉}

围绕“MESI 直觉”，先把它放回一个具体问题：只读表和热点计数器为什么都是共享却表现不同。如果这个问题只在真实线上服务里出现，读者很容易把它当成局部技巧；但在多线程计数和大数组归约中，它会沿着输入、执行、同步、I/O 和提交一路传播。所以讲解时不能只写定义，而要让读者看到它怎样从一个小失败变成系统性约束。

朴素写法通常是“把共享数据一概视为危险”。这句话看似简单，却包含几个默认前提：状态只有一个拥有者，执行顺序可预测，资源足够近，失败不会穿过边界，观察结果能够代表真实成本。这些前提在教学小程序中可能成立，在多线程计数和大数组归约里却会被输入规模、线程交错、硬件拓扑或故障注入逐个打破。

失败信号是“读共享可复制，高频写共享要反复失效或迁移”。注意，这个失败不一定表现为崩溃。它也可能表现为吞吐不再增长、p99 抖动、重启后结果多了一份、某个分区长尾、CPU 使用率很高但有效工作很少。教材要把这些信号和机制绑定起来，否则读者只会背术语，遇到真实现象时仍然不知道从哪里下手。

更可靠的推导顺序是先写出不变量。对MESI 直觉来说，不变量不是一句口号，而是本章要保护的事实：哪些数据只能写一次，哪些状态可以重试，哪些结果可以被缓存，哪些成本必须摊销，哪些错误必须外显。只要不变量没有写清，后面的代码、实验和优化都没有稳定坐标。

然后才引入模型：Modified、Exclusive、Shared、Invalid 可以帮助理解状态迁移。这个模型应同时解释正确性和成本。正确性部分回答“为什么结果可信”，成本部分回答“为什么这个实现会快或慢”。如果一个模型只能解释速度，却解释不了失败恢复，它不适合系统教材；如果只能解释语义，却完全不关心 cache、调度、I/O 或网络成本，也不适合第二册。

实验应围绕“只改变一个变量”设计。针对MESI 直觉，最小实验可以保留朴素版本，再实现一个改进版本，输入规模从小到大，线程数或并发度从一到目标机器上限，错误注入从无到有。每一组实验都应记录 reference 结果、核心指标、环境和命令。这样读者看到差异时，可以判断差异来自机制本身，而不是来自初始化、缓存冷热、调度迁移或文件系统状态。

观察手段是：用文本状态图解释读多写少和写热点。观察不是为了堆工具名，而是为了回答一个具体假设。若假设是共享写导致扩展性下降，就应看线程数曲线、写入布局和 cache 相关指标；若假设是提交边界错误，就应做崩溃点矩阵；若假设是队列背压，就应看水位、等待和上游速率。工具输出必须回到假设，不要把截图或命令输出当成结论。

最后要讲边界：真实协议更复杂，教学模型只用于解释软件现象。高质量教材必须把反例放在正文里。读者需要知道什么时候该用这个机制，什么时候它会制造新问题，什么时候应该退回更简单方案。比如一个单线程工具没有并发共享，强行引入复杂运行时只会增加维护成本；一个没有恢复要求的临时分析脚本，也不必承担完整 manifest 协议。

把这一点落实到代码审查，可以要求每个涉及MESI 直觉的改动都回答五个问题：朴素版本是什么，新增机制保护了哪个不变量，新增成本在哪里，如何回退，哪些测试证明边界。这五个问题比“用了某某技术”更能判断质量，也能防止团队在没有证据时追逐复杂实现。

\topic{深挖：Store buffer}

围绕“Store buffer”，先把它放回一个具体问题：写 data 再写 flag 为什么仍需要同步语义。如果这个问题只在单线程小输入里出现，读者很容易把它当成局部技巧；但在多线程计数和大数组归约中，它会沿着输入、执行、同步、I/O 和提交一路传播。所以讲解时不能只写定义，而要让读者看到它怎样从一个小失败变成系统性约束。

朴素写法通常是“相信源码顺序就是其他线程看到的顺序”。这句话看似简单，却包含几个默认前提：状态只有一个拥有者，执行顺序可预测，资源足够近，失败不会穿过边界，观察结果能够代表真实成本。这些前提在教学小程序中可能成立，在多线程计数和大数组归约里却会被输入规模、线程交错、硬件拓扑或故障注入逐个打破。

失败信号是“写入可能先在缓冲中，对其他核心的可见顺序必须由同步建立”。注意，这个失败不一定表现为崩溃。它也可能表现为吞吐不再增长、p99 抖动、重启后结果多了一份、某个分区长尾、CPU 使用率很高但有效工作很少。教材要把这些信号和机制绑定起来，否则读者只会背术语，遇到真实现象时仍然不知道从哪里下手。

更可靠的推导顺序是先写出不变量。对Store buffer来说，不变量不是一句口号，而是本章要保护的事实：哪些数据只能写一次，哪些状态可以重试，哪些结果可以被缓存，哪些成本必须摊销，哪些错误必须外显。只要不变量没有写清，后面的代码、实验和优化都没有稳定坐标。

然后才引入模型：store buffer 和内存模型共同决定发布可见性。这个模型应同时解释正确性和成本。正确性部分回答“为什么结果可信”，成本部分回答“为什么这个实现会快或慢”。如果一个模型只能解释速度，却解释不了失败恢复，它不适合系统教材；如果只能解释语义，却完全不关心 cache、调度、I/O 或网络成本，也不适合第二册。

实验应围绕“只改变一个变量”设计。针对Store buffer，最小实验可以保留朴素版本，再实现一个改进版本，输入规模从小到大，线程数或并发度从一到目标机器上限，错误注入从无到有。每一组实验都应记录 reference 结果、核心指标、环境和命令。这样读者看到差异时，可以判断差异来自机制本身，而不是来自初始化、缓存冷热、调度迁移或文件系统状态。

观察手段是：用 release acquire 和错误普通变量版本做对照。观察不是为了堆工具名，而是为了回答一个具体假设。若假设是共享写导致扩展性下降，就应看线程数曲线、写入布局和 cache 相关指标；若假设是提交边界错误，就应做崩溃点矩阵；若假设是队列背压，就应看水位、等待和上游速率。工具输出必须回到假设，不要把截图或命令输出当成结论。

最后要讲边界：不要用硬件一致性为 C++ 数据竞争辩护。高质量教材必须把反例放在正文里。读者需要知道什么时候该用这个机制，什么时候它会制造新问题，什么时候应该退回更简单方案。比如一个单线程工具没有并发共享，强行引入复杂运行时只会增加维护成本；一个没有恢复要求的临时分析脚本，也不必承担完整 manifest 协议。

把这一点落实到代码审查，可以要求每个涉及Store buffer的改动都回答五个问题：朴素版本是什么，新增机制保护了哪个不变量，新增成本在哪里，如何回退，哪些测试证明边界。这五个问题比“用了某某技术”更能判断质量，也能防止团队在没有证据时追逐复杂实现。

\topic{深挖：NUMA first touch}

围绕“NUMA first touch”，先把它放回一个具体问题：为什么同一数组在不同初始化方式下多线程带宽不同。如果这个问题只在单机多线程里出现，读者很容易把它当成局部技巧；但在多线程计数和大数组归约中，它会沿着输入、执行、同步、I/O 和提交一路传播。所以讲解时不能只写定义，而要让读者看到它怎样从一个小失败变成系统性约束。

朴素写法通常是“主线程分配并初始化所有页面”。这句话看似简单，却包含几个默认前提：状态只有一个拥有者，执行顺序可预测，资源足够近，失败不会穿过边界，观察结果能够代表真实成本。这些前提在教学小程序中可能成立，在多线程计数和大数组归约里却会被输入规模、线程交错、硬件拓扑或故障注入逐个打破。

失败信号是“页面可能集中在一个节点，其他节点 worker 远端访问”。注意，这个失败不一定表现为崩溃。它也可能表现为吞吐不再增长、p99 抖动、重启后结果多了一份、某个分区长尾、CPU 使用率很高但有效工作很少。教材要把这些信号和机制绑定起来，否则读者只会背术语，遇到真实现象时仍然不知道从哪里下手。

更可靠的推导顺序是先写出不变量。对NUMA first touch来说，不变量不是一句口号，而是本章要保护的事实：哪些数据只能写一次，哪些状态可以重试，哪些结果可以被缓存，哪些成本必须摊销，哪些错误必须外显。只要不变量没有写清，后面的代码、实验和优化都没有稳定坐标。

然后才引入模型：NUMA 把内存位置引入成本模型，first touch 影响页面放置。这个模型应同时解释正确性和成本。正确性部分回答“为什么结果可信”，成本部分回答“为什么这个实现会快或慢”。如果一个模型只能解释速度，却解释不了失败恢复，它不适合系统教材；如果只能解释语义，却完全不关心 cache、调度、I/O 或网络成本，也不适合第二册。

实验应围绕“只改变一个变量”设计。针对NUMA first touch，最小实验可以保留朴素版本，再实现一个改进版本，输入规模从小到大，线程数或并发度从一到目标机器上限，错误注入从无到有。每一组实验都应记录 reference 结果、核心指标、环境和命令。这样读者看到差异时，可以判断差异来自机制本身，而不是来自初始化、缓存冷热、调度迁移或文件系统状态。

观察手段是：比较绑定与未绑定、single touch 与 parallel touch。观察不是为了堆工具名，而是为了回答一个具体假设。若假设是共享写导致扩展性下降，就应看线程数曲线、写入布局和 cache 相关指标；若假设是提交边界错误，就应做崩溃点矩阵；若假设是队列背压，就应看水位、等待和上游速率。工具输出必须回到假设，不要把截图或命令输出当成结论。

最后要讲边界：没有 NUMA 机器时不能声称验证了远端内存。高质量教材必须把反例放在正文里。读者需要知道什么时候该用这个机制，什么时候它会制造新问题，什么时候应该退回更简单方案。比如一个单线程工具没有并发共享，强行引入复杂运行时只会增加维护成本；一个没有恢复要求的临时分析脚本，也不必承担完整 manifest 协议。

把这一点落实到代码审查，可以要求每个涉及NUMA first touch的改动都回答五个问题：朴素版本是什么，新增机制保护了哪个不变量，新增成本在哪里，如何回退，哪些测试证明边界。这五个问题比“用了某某技术”更能判断质量，也能防止团队在没有证据时追逐复杂实现。

\topic{深挖：互连饱和}

围绕“互连饱和”，先把它放回一个具体问题：每个线程单独跑都快，一起跑为什么都慢。如果这个问题只在NUMA 或异构核心里出现，读者很容易把它当成局部技巧；但在多线程计数和大数组归约中，它会沿着输入、执行、同步、I/O 和提交一路传播。所以讲解时不能只写定义，而要让读者看到它怎样从一个小失败变成系统性约束。

朴素写法通常是“继续增加线程寻找吞吐”。这句话看似简单，却包含几个默认前提：状态只有一个拥有者，执行顺序可预测，资源足够近，失败不会穿过边界，观察结果能够代表真实成本。这些前提在教学小程序中可能成立，在多线程计数和大数组归约里却会被输入规模、线程交错、硬件拓扑或故障注入逐个打破。

失败信号是“内存控制器、LLC 或 socket 互连达到上限”。注意，这个失败不一定表现为崩溃。它也可能表现为吞吐不再增长、p99 抖动、重启后结果多了一份、某个分区长尾、CPU 使用率很高但有效工作很少。教材要把这些信号和机制绑定起来，否则读者只会背术语，遇到真实现象时仍然不知道从哪里下手。

更可靠的推导顺序是先写出不变量。对互连饱和来说，不变量不是一句口号，而是本章要保护的事实：哪些数据只能写一次，哪些状态可以重试，哪些结果可以被缓存，哪些成本必须摊销，哪些错误必须外显。只要不变量没有写清，后面的代码、实验和优化都没有稳定坐标。

然后才引入模型：核心、缓存、内存控制器和 socket 之间共享有限互连资源。这个模型应同时解释正确性和成本。正确性部分回答“为什么结果可信”，成本部分回答“为什么这个实现会快或慢”。如果一个模型只能解释速度，却解释不了失败恢复，它不适合系统教材；如果只能解释语义，却完全不关心 cache、调度、I/O 或网络成本，也不适合第二册。

实验应围绕“只改变一个变量”设计。针对互连饱和，最小实验可以保留朴素版本，再实现一个改进版本，输入规模从小到大，线程数或并发度从一到目标机器上限，错误注入从无到有。每一组实验都应记录 reference 结果、核心指标、环境和命令。这样读者看到差异时，可以判断差异来自机制本身，而不是来自初始化、缓存冷热、调度迁移或文件系统状态。

观察手段是：画线程数到吞吐曲线，并记录带宽和 LLC 指标。观察不是为了堆工具名，而是为了回答一个具体假设。若假设是共享写导致扩展性下降，就应看线程数曲线、写入布局和 cache 相关指标；若假设是提交边界错误，就应做崩溃点矩阵；若假设是队列背压，就应看水位、等待和上游速率。工具输出必须回到假设，不要把截图或命令输出当成结论。

最后要讲边界：带宽瓶颈下更多线程可能只增加争用。高质量教材必须把反例放在正文里。读者需要知道什么时候该用这个机制，什么时候它会制造新问题，什么时候应该退回更简单方案。比如一个单线程工具没有并发共享，强行引入复杂运行时只会增加维护成本；一个没有恢复要求的临时分析脚本，也不必承担完整 manifest 协议。

把这一点落实到代码审查，可以要求每个涉及互连饱和的改动都回答五个问题：朴素版本是什么，新增机制保护了哪个不变量，新增成本在哪里，如何回退，哪些测试证明边界。这五个问题比“用了某某技术”更能判断质量，也能防止团队在没有证据时追逐复杂实现。

\topic{深挖：分层归约}

围绕“分层归约”，先把它放回一个具体问题：为什么先本地合并再跨节点合并比所有线程直接写全局结果更稳。如果这个问题只在本机分布式模拟里出现，读者很容易把它当成局部技巧；但在多线程计数和大数组归约中，它会沿着输入、执行、同步、I/O 和提交一路传播。所以讲解时不能只写定义，而要让读者看到它怎样从一个小失败变成系统性约束。

朴素写法通常是“最后一个全局变量收集所有 partial”。这句话看似简单，却包含几个默认前提：状态只有一个拥有者，执行顺序可预测，资源足够近，失败不会穿过边界，观察结果能够代表真实成本。这些前提在教学小程序中可能成立，在多线程计数和大数组归约里却会被输入规模、线程交错、硬件拓扑或故障注入逐个打破。

失败信号是“跨节点高频写或集中合并会制造热点”。注意，这个失败不一定表现为崩溃。它也可能表现为吞吐不再增长、p99 抖动、重启后结果多了一份、某个分区长尾、CPU 使用率很高但有效工作很少。教材要把这些信号和机制绑定起来，否则读者只会背术语，遇到真实现象时仍然不知道从哪里下手。

更可靠的推导顺序是先写出不变量。对分层归约来说，不变量不是一句口号，而是本章要保护的事实：哪些数据只能写一次，哪些状态可以重试，哪些结果可以被缓存，哪些成本必须摊销，哪些错误必须外显。只要不变量没有写清，后面的代码、实验和优化都没有稳定坐标。

然后才引入模型：分层归约把高频路径留在近处，把远端通信压缩到少量结果。这个模型应同时解释正确性和成本。正确性部分回答“为什么结果可信”，成本部分回答“为什么这个实现会快或慢”。如果一个模型只能解释速度，却解释不了失败恢复，它不适合系统教材；如果只能解释语义，却完全不关心 cache、调度、I/O 或网络成本，也不适合第二册。

实验应围绕“只改变一个变量”设计。针对分层归约，最小实验可以保留朴素版本，再实现一个改进版本，输入规模从小到大，线程数或并发度从一到目标机器上限，错误注入从无到有。每一组实验都应记录 reference 结果、核心指标、环境和命令。这样读者看到差异时，可以判断差异来自机制本身，而不是来自初始化、缓存冷热、调度迁移或文件系统状态。

观察手段是：比较 worker、NUMA node、global 三层耗时。观察不是为了堆工具名，而是为了回答一个具体假设。若假设是共享写导致扩展性下降，就应看线程数曲线、写入布局和 cache 相关指标；若假设是提交边界错误，就应做崩溃点矩阵；若假设是队列背压，就应看水位、等待和上游速率。工具输出必须回到假设，不要把截图或命令输出当成结论。

最后要讲边界：层级太多会增加阶段开销，小数据未必值得。高质量教材必须把反例放在正文里。读者需要知道什么时候该用这个机制，什么时候它会制造新问题，什么时候应该退回更简单方案。比如一个单线程工具没有并发共享，强行引入复杂运行时只会增加维护成本；一个没有恢复要求的临时分析脚本，也不必承担完整 manifest 协议。

把这一点落实到代码审查，可以要求每个涉及分层归约的改动都回答五个问题：朴素版本是什么，新增机制保护了哪个不变量，新增成本在哪里，如何回退，哪些测试证明边界。这五个问题比“用了某某技术”更能判断质量，也能防止团队在没有证据时追逐复杂实现。

\topic{深挖：拓扑感知队列}

围绕“拓扑感知队列”，先把它放回一个具体问题：全局任务队列为什么在多 socket 上成为热点。如果这个问题只在真实线上服务里出现，读者很容易把它当成局部技巧；但在多线程计数和大数组归约中，它会沿着输入、执行、同步、I/O 和提交一路传播。所以讲解时不能只写定义，而要让读者看到它怎样从一个小失败变成系统性约束。

朴素写法通常是“所有 worker 共享一个 MPMC 队列”。这句话看似简单，却包含几个默认前提：状态只有一个拥有者，执行顺序可预测，资源足够近，失败不会穿过边界，观察结果能够代表真实成本。这些前提在教学小程序中可能成立，在多线程计数和大数组归约里却会被输入规模、线程交错、硬件拓扑或故障注入逐个打破。

失败信号是“head、tail、锁和槽位状态在节点间迁移”。注意，这个失败不一定表现为崩溃。它也可能表现为吞吐不再增长、p99 抖动、重启后结果多了一份、某个分区长尾、CPU 使用率很高但有效工作很少。教材要把这些信号和机制绑定起来，否则读者只会背术语，遇到真实现象时仍然不知道从哪里下手。

更可靠的推导顺序是先写出不变量。对拓扑感知队列来说，不变量不是一句口号，而是本章要保护的事实：哪些数据只能写一次，哪些状态可以重试，哪些结果可以被缓存，哪些成本必须摊销，哪些错误必须外显。只要不变量没有写清，后面的代码、实验和优化都没有稳定坐标。

然后才引入模型：本地队列和工作窃取把高频访问局部化，低频偷取承担远端成本。这个模型应同时解释正确性和成本。正确性部分回答“为什么结果可信”，成本部分回答“为什么这个实现会快或慢”。如果一个模型只能解释速度，却解释不了失败恢复，它不适合系统教材；如果只能解释语义，却完全不关心 cache、调度、I/O 或网络成本，也不适合第二册。

实验应围绕“只改变一个变量”设计。针对拓扑感知队列，最小实验可以保留朴素版本，再实现一个改进版本，输入规模从小到大，线程数或并发度从一到目标机器上限，错误注入从无到有。每一组实验都应记录 reference 结果、核心指标、环境和命令。这样读者看到差异时，可以判断差异来自机制本身，而不是来自初始化、缓存冷热、调度迁移或文件系统状态。

观察手段是：记录本地 pop、同节点 steal、跨节点 steal 次数。观察不是为了堆工具名，而是为了回答一个具体假设。若假设是共享写导致扩展性下降，就应看线程数曲线、写入布局和 cache 相关指标；若假设是提交边界错误，就应做崩溃点矩阵；若假设是队列背压，就应看水位、等待和上游速率。工具输出必须回到假设，不要把截图或命令输出当成结论。

最后要讲边界：队列分片会增加负载均衡和关闭语义复杂度。高质量教材必须把反例放在正文里。读者需要知道什么时候该用这个机制，什么时候它会制造新问题，什么时候应该退回更简单方案。比如一个单线程工具没有并发共享，强行引入复杂运行时只会增加维护成本；一个没有恢复要求的临时分析脚本，也不必承担完整 manifest 协议。

把这一点落实到代码审查，可以要求每个涉及拓扑感知队列的改动都回答五个问题：朴素版本是什么，新增机制保护了哪个不变量，新增成本在哪里，如何回退，哪些测试证明边界。这五个问题比“用了某某技术”更能判断质量，也能防止团队在没有证据时追逐复杂实现。

\topic{深挖：拓扑记录}

围绕“拓扑记录”，先把它放回一个具体问题：为什么只写 worker count 不足以复现实验。如果这个问题只在单线程小输入里出现，读者很容易把它当成局部技巧；但在多线程计数和大数组归约中，它会沿着输入、执行、同步、I/O 和提交一路传播。所以讲解时不能只写定义，而要让读者看到它怎样从一个小失败变成系统性约束。

朴素写法通常是“报告只记录线程数”。这句话看似简单，却包含几个默认前提：状态只有一个拥有者，执行顺序可预测，资源足够近，失败不会穿过边界，观察结果能够代表真实成本。这些前提在教学小程序中可能成立，在多线程计数和大数组归约里却会被输入规模、线程交错、硬件拓扑或故障注入逐个打破。

失败信号是“SMT、socket、NUMA node、频率和容器限制都会改变结果”。注意，这个失败不一定表现为崩溃。它也可能表现为吞吐不再增长、p99 抖动、重启后结果多了一份、某个分区长尾、CPU 使用率很高但有效工作很少。教材要把这些信号和机制绑定起来，否则读者只会背术语，遇到真实现象时仍然不知道从哪里下手。

更可靠的推导顺序是先写出不变量。对拓扑记录来说，不变量不是一句口号，而是本章要保护的事实：哪些数据只能写一次，哪些状态可以重试，哪些结果可以被缓存，哪些成本必须摊销，哪些错误必须外显。只要不变量没有写清，后面的代码、实验和优化都没有稳定坐标。

然后才引入模型：拓扑记录把 CPU id、core、node、cache 共享和绑定策略写进报告。这个模型应同时解释正确性和成本。正确性部分回答“为什么结果可信”，成本部分回答“为什么这个实现会快或慢”。如果一个模型只能解释速度，却解释不了失败恢复，它不适合系统教材；如果只能解释语义，却完全不关心 cache、调度、I/O 或网络成本，也不适合第二册。

实验应围绕“只改变一个变量”设计。针对拓扑记录，最小实验可以保留朴素版本，再实现一个改进版本，输入规模从小到大，线程数或并发度从一到目标机器上限，错误注入从无到有。每一组实验都应记录 reference 结果、核心指标、环境和命令。这样读者看到差异时，可以判断差异来自机制本身，而不是来自初始化、缓存冷热、调度迁移或文件系统状态。

观察手段是：保存 lscpu、numactl、cpuset 和程序绑定配置。观察不是为了堆工具名，而是为了回答一个具体假设。若假设是共享写导致扩展性下降，就应看线程数曲线、写入布局和 cache 相关指标；若假设是提交边界错误，就应做崩溃点矩阵；若假设是队列背压，就应看水位、等待和上游速率。工具输出必须回到假设，不要把截图或命令输出当成结论。

最后要讲边界：生产系统是否绑定要看部署，不是所有服务都适合强绑定。高质量教材必须把反例放在正文里。读者需要知道什么时候该用这个机制，什么时候它会制造新问题，什么时候应该退回更简单方案。比如一个单线程工具没有并发共享，强行引入复杂运行时只会增加维护成本；一个没有恢复要求的临时分析脚本，也不必承担完整 manifest 协议。

把这一点落实到代码审查，可以要求每个涉及拓扑记录的改动都回答五个问题：朴素版本是什么，新增机制保护了哪个不变量，新增成本在哪里，如何回退，哪些测试证明边界。这五个问题比“用了某某技术”更能判断质量，也能防止团队在没有证据时追逐复杂实现。

\topic{案例深化：从朴素版到工程版}

案例深化“全局 atomic 计数器”要按三次迭代写。第一次只写 reference：多个线程每条记录都更新同一个统计值。reference 的代码可以慢，但必须把输入、输出、错误和边界写清。第二次写朴素工程版本：使用 relaxed fetch add。这一步不能省，因为读者需要看到直觉方案为什么诱人，也需要有一个失败对象。

第三次才写改进版本：改为 per worker 或 per NUMA shard。改进说明要避免“因为更高级所以更快”这种表达，而要讲清它改变了哪条路径：减少了数据移动、减少了共享写、减少了系统调用、减少了重试放大，还是把不可恢复状态变成可恢复状态。

验收方式是：比较扩展曲线和每次操作平均成本。验收报告至少包含一张对照表，一列是朴素版本，一列是改进版本，一列是反例。反例很重要，因为它能告诉读者改进不是什么时候都正确。如果一个案例没有反例，往往说明分析还停在宣传层面。
案例深化“双 socket 数组归约”要按三次迭代写。第一次只写 reference：数组大小超过 LLC，worker 分布在两个 socket。reference 的代码可以慢，但必须把输入、输出、错误和边界写清。第二次写朴素工程版本：主线程初始化全部页面。这一步不能省，因为读者需要看到直觉方案为什么诱人，也需要有一个失败对象。

第三次才写改进版本：按 worker 绑定和 first touch 分块初始化。改进说明要避免“因为更高级所以更快”这种表达，而要讲清它改变了哪条路径：减少了数据移动、减少了共享写、减少了系统调用、减少了重试放大，还是把不可恢复状态变成可恢复状态。

验收方式是：观察带宽、远端访问迹象和节点内合并成本。验收报告至少包含一张对照表，一列是朴素版本，一列是改进版本，一列是反例。反例很重要，因为它能告诉读者改进不是什么时候都正确。如果一个案例没有反例，往往说明分析还停在宣传层面。
案例深化“全局队列与本地队列”要按三次迭代写。第一次只写 reference：大量短任务由所有 worker 抢占执行。reference 的代码可以慢，但必须把输入、输出、错误和边界写清。第二次写朴素工程版本：一个全局队列存所有任务。这一步不能省，因为读者需要看到直觉方案为什么诱人，也需要有一个失败对象。

第三次才写改进版本：每节点队列加窃取。改进说明要避免“因为更高级所以更快”这种表达，而要讲清它改变了哪条路径：减少了数据移动、减少了共享写、减少了系统调用、减少了重试放大，还是把不可恢复状态变成可恢复状态。

验收方式是：比较短任务和长任务两种粒度下的队列成本。验收报告至少包含一张对照表，一列是朴素版本，一列是改进版本，一列是反例。反例很重要，因为它能告诉读者改进不是什么时候都正确。如果一个案例没有反例，往往说明分析还停在宣传层面。

\topic{实验矩阵：让结论可复现}

围绕一致性所有权，实验矩阵至少包含正常输入、边界输入、压力输入和错误输入四类。正常输入验证功能，边界输入验证不变量，压力输入验证成本模型，错误输入验证恢复和报告。其中压力输入不只是“大”，还要能触发本章关心的形状：随机访问、热点 key、短任务、慢 I/O、重复消息或跨节点访问。

围绕MESI 直觉，实验矩阵至少包含正常输入、边界输入、压力输入和错误输入四类。正常输入验证功能，边界输入验证不变量，压力输入验证成本模型，错误输入验证恢复和报告。其中压力输入不只是“大”，还要能触发本章关心的形状：随机访问、热点 key、短任务、慢 I/O、重复消息或跨节点访问。

围绕Store buffer，实验矩阵至少包含正常输入、边界输入、压力输入和错误输入四类。正常输入验证功能，边界输入验证不变量，压力输入验证成本模型，错误输入验证恢复和报告。其中压力输入不只是“大”，还要能触发本章关心的形状：随机访问、热点 key、短任务、慢 I/O、重复消息或跨节点访问。

围绕NUMA first touch，实验矩阵至少包含正常输入、边界输入、压力输入和错误输入四类。正常输入验证功能，边界输入验证不变量，压力输入验证成本模型，错误输入验证恢复和报告。其中压力输入不只是“大”，还要能触发本章关心的形状：随机访问、热点 key、短任务、慢 I/O、重复消息或跨节点访问。

围绕互连饱和，实验矩阵至少包含正常输入、边界输入、压力输入和错误输入四类。正常输入验证功能，边界输入验证不变量，压力输入验证成本模型，错误输入验证恢复和报告。其中压力输入不只是“大”，还要能触发本章关心的形状：随机访问、热点 key、短任务、慢 I/O、重复消息或跨节点访问。

围绕分层归约，实验矩阵至少包含正常输入、边界输入、压力输入和错误输入四类。正常输入验证功能，边界输入验证不变量，压力输入验证成本模型，错误输入验证恢复和报告。其中压力输入不只是“大”，还要能触发本章关心的形状：随机访问、热点 key、短任务、慢 I/O、重复消息或跨节点访问。

围绕拓扑感知队列，实验矩阵至少包含正常输入、边界输入、压力输入和错误输入四类。正常输入验证功能，边界输入验证不变量，压力输入验证成本模型，错误输入验证恢复和报告。其中压力输入不只是“大”，还要能触发本章关心的形状：随机访问、热点 key、短任务、慢 I/O、重复消息或跨节点访问。

围绕拓扑记录，实验矩阵至少包含正常输入、边界输入、压力输入和错误输入四类。正常输入验证功能，边界输入验证不变量，压力输入验证成本模型，错误输入验证恢复和报告。其中压力输入不只是“大”，还要能触发本章关心的形状：随机访问、热点 key、短任务、慢 I/O、重复消息或跨节点访问。

\topic{Linux 源码阅读深化}

本章的 Linux 阅读入口仍然保持窄范围：\filepath{kernel/sched/topology.c}、\filepath{mm/mempolicy.c}、\filepath{kernel/sched/fair.c}、\filepath{kernel/locking/qspinlock.c}。阅读时不要追求覆盖率，而要追求因果链。从一个用户态现象出发，找到内核中的状态对象，再看状态怎样被修改、等待怎样被挂起、唤醒怎样发生、错误怎样返回。

源码阅读可以按四栏笔记整理。第一栏写用户态操作，例如线程等待、缺页、写文件、发送消息或计时。第二栏写内核对象，例如 task、VMA、page、inode、wait queue、bio、socket buffer。第三栏写关键状态变化，例如 runnable 到 sleeping、clean page 到 dirty page、pending task 到 committed task。第四栏写可观察指标或实验，例如 context switch、page fault、writeback、queue depth、retry count。

不要把 Linux 源码当成术语来源。源码阅读的目标是训练映射能力：把书中的抽象边界映射到真实系统里的结构和状态。读者不需要一次读懂整个内核，但要能解释一个最小现象为什么发生、哪些路径参与、哪些结论受当前内核版本和硬件限制。

\topic{项目验收：把本章能力写进仓库}

贯穿项目的验收不能只说“功能跑通”。本章至少要留下一个可重复实验、一个失败注入、一个指标输出和一个设计说明。实验说明机制是否真的被触发；失败注入说明边界是否真的被处理；指标输出说明结论是否有证据；设计说明让后续章节能复用这个模块。

本章项目检查点可以具体化为：1. 为计数器实现全局和分片两种路径
2. 记录 CPU 拓扑和线程绑定
3. 加入分层归约实验
4. 统计本地队列和跨节点窃取次数
5. 写 NUMA 可验证与不可验证边界。这些检查点要进入仓库，而不是停留在阅读笔记。如果某个检查点因为当前设备权限、硬件或系统 API 不支持而无法完成，也要写清降级方案和未验证风险。

项目验收还应要求读者写一段复盘：本章新增机制让系统多了什么能力，也让系统多了什么复杂度。比如加入背压后内存更稳定，但生产者会等待；加入 route epoch 后旧请求能被拒绝，但所有消息都必须携带版本；加入 SIMD 后热内核更快，但尾部和数值语义要额外测试。这种复盘能把知识从“会用”推进到“会判断”。



\topic{最终收束：本章应形成的判断力}

读完本章，读者不应只记住一致性所有权、拓扑记录这些词，而应能围绕多线程计数和大数组归约做一次完整判断。完整判断从问题开始：现象是什么，朴素方案是什么，失败信号是什么，保护的不变量是什么，成本从哪里移动到哪里。若无法回答这些问题，就说明还停留在概念层，没有真正进入系统层。

本章最重要的反例是：线程数量增加后，瓶颈从计算转移到 cache line 所有权、互连带宽和远端内存。遇到类似反例时，第一步不是换技术，而是缩小问题。先固定输入，固定环境，保留 reference，再一次只改变一个变量。如果改动同时改变布局、线程数、批大小、同步方式和提交策略，即使结果变快，也无法知道原因。高质量工程实验必须克制变量。

以案例“全局 atomic 计数器”为例，最终报告应包含三段。第一段写多个线程每条记录都更新同一个统计值，说明读者要解决的不是抽象题，而是一个有输入、有状态、有失败方式的任务。第二段写朴素版本：使用 relaxed fetch add，并诚实说明它为什么吸引人。第三段写改进版本：改为 per worker 或 per NUMA shard，再用比较扩展曲线和每次操作平均成本验收。报告中若没有朴素版本，读者会误以为最终设计是凭空出现；若没有反例，读者会误以为最终设计永远正确。

本章还应留下一个审查表。正确性方面，检查输入合同、状态转移、所有权、重复执行、关闭和恢复。性能方面，检查数据移动、共享写、排队、系统调用、批大小、缓存冷热和拓扑。可观测性方面，检查是否有阶段耗时、队列水位、错误分类、任务身份和原始样本。每一项都要能落到代码、测试或报告，不要只停在口头承诺。

贯穿项目里，本章至少要完成这些检查点：为计数器实现全局和分片两种路径、记录 CPU 拓扑和线程绑定、加入分层归约实验、统计本地队列和跨节点窃取次数、写 NUMA 可验证与不可验证边界。完成后不要急着进入下一章，要先写一页复盘：本章新增能力解决了什么失败，新增了哪些复杂度，哪些结论只在当前机器上验证，哪些结论可以迁移到更大系统。这种复盘会让整本书的知识保持连续，不会变成每章讲完就散掉的材料。

最后，用一句话收束本章：系统能力来自清楚的边界、可验证的不变量和可复现的证据。无论本章讨论的是硬件执行、内存层级、同步、运行时、I/O 还是分布式协议，只要缺少边界、不变量和证据，复杂实现就会变成风险来源。反过来，只要这三件事稳住，读者就能把一个朴素程序逐步推进成可靠的计算系统。




\topic{过线补充：常见误区、验收题和迁移能力}

本章最容易出现的误区，是把一致性所有权、MESI 直觉、Store buffer、NUMA first touch当成互相独立的知识点。在多线程计数和大数组归约中，它们其实围绕同一个失败展开：线程数量增加后，瓶颈从计算转移到 cache line 所有权、互连带宽和远端内存。如果读者只能分别解释这些概念，却不能说明它们在同一条数据流里怎样互相影响，就还没有真正掌握本章。例如一个优化减少了计算时间，却增加了队列等待；一个同步方案修复了正确性，却制造了共享写热点；一个提交协议提高了可靠性，却增加了持久化延迟。这些都不是矛盾，而是系统设计中的成本迁移。

本章的验收题可以这样设计：先给出一个最小版本的多线程计数和大数组归约，要求读者写出 reference；再引入一个压力条件，要求读者指出朴素方案会在哪里失败；最后要求读者选择本章机制中的一项或两项做改进。答案不能只给最终代码，还要给不变量、实验矩阵和反例。如果某个答案没有说明失败后如何恢复、如何观察、如何判断优化是否值得，就不能算完整答案。

围绕案例全局 atomic 计数器、双 socket 数组归约、全局队列与本地队列，报告还应补一个迁移问题：同样方法迁移到更大输入、更多线程、不同硬件或本机分布式模拟时，哪些结论保持，哪些结论要重新测。保持的通常是语义边界和不变量；需要重测的通常是吞吐、延迟、批大小、缓存效果、锁竞争和 I/O 行为。这能训练读者区分“原理可以迁移”和“数字不能照搬”。

最后再强调一次本书的写法要求：先从问题出发，再推导机制，再落到实验。不要把实验放成装饰，也不要把 Linux 源码阅读放成资料链接。实验必须检验一个假设，源码阅读必须解释一个用户态现象，项目代码必须保护一个明确边界。读者能按这个顺序复述本章，才说明本章内容真正连贯起来。




\topic{事故复盘式走查}

为了把本章内容真正串起来，可以把多线程计数和大数组归约写成一次小型事故复盘。事故现象不是一句“系统慢了”或“结果错了”，而要具体到可观察事实：哪一批输入、哪个阶段、哪一个指标、哪一种失败注入触发了问题。如果现象描述不具体，后面的一致性所有权、MESI 直觉和Store buffer就只能变成猜测。

复盘第一步是还原朴素设计。写清当时为什么选择它，它解决了什么简单问题，又默认了哪些条件。很多坏设计一开始并不坏，只是从小输入、小并发、无失败的条件迁移到多线程计数和大数组归约时，没有同步升级边界。这也是本册反复强调从朴素方案出发的原因：只有理解朴素方案为什么成立，才能准确说明它为什么在新条件下失败。

复盘第二步是列出候选假设，但每个假设都必须有可证伪实验。如果怀疑一致性所有权相关路径，就构造只改变这一因素的对照；如果怀疑MESI 直觉，就记录它对应的状态或指标；如果怀疑Store buffer，就找出它在代码里的所有入口和退出点。不能把所有怀疑同时改掉。一次修改多个因素会让复盘变成经验叙事，而不是工程证据。

复盘第三步是修复。修复不等于把代码改到通过测试，而是把不变量写回系统。本章的不变量来自“线程数量增加后，瓶颈从计算转移到 cache line 所有权、互连带宽和远端内存”这个失败主题。因此修复说明要写清：新增字段保护什么，新增队列容量限制什么，新增版本号拒绝什么，新增 reference 比较什么，新增指标暴露什么。如果修复无法对应到不变量，就很可能只是局部补丁。

复盘第四步是验收。以“全局队列与本地队列”为例，验收不应只跑正常输入，还要跑边界输入、压力输入和错误输入。正常输入证明功能还在，边界输入证明不变量没有被破坏，压力输入证明成本模型仍然成立，错误输入证明失败不会越过提交边界。验收报告要保留原始命令和数据，否则下一次回归无法判断问题是否真正修复。

最后要写“未解决问题”。例如拓扑记录相关边界可能还没有在当前设备上完全验证，某些 Linux 计数器可能因为权限不可用，某些 NUMA 结论可能因为硬件不支持只能停留在设计层。把未验证风险写出来不是降低质量，而是防止教材把假设写成事实。高质量书籍要敢于说明证据边界，这比把每个结论都写成绝对经验更可靠。

读者完成本章后，可以用这份复盘模板检查自己的学习结果：能否描述事故，能否还原朴素设计，能否提出可证伪假设，能否把修复绑定到不变量，能否设计验收矩阵，能否写出未验证风险。如果六项都能完成，本章知识就已经从概念记忆转成工程判断。如果只能完成其中一两项，就应回到本章前面的案例重新走一遍，而不是急着进入下一章。



\section{本章落到贯穿项目}

Compute Systems Engine 在本章应加入三个设计点。第一，拓扑记录：benchmark 输出可见 CPU 数、用户指定 worker 数、是否启用绑定、是否能读取 NUMA 信息。第二，分层 partial：即使当前只实现每 worker partial，也在设计文档中说明未来可以扩展成每 NUMA 节点 partial。第三，队列分片路线：线程池初版可以用全局队列，但要记录全局队列可能成为共享热点，后续工作窃取章节再拆成本地队列。

本章的验收不是一定在手机或单 socket 机器上测出 NUMA 差异，而是实验设计正确。没有 NUMA 设备时，报告应写明限制：可以验证 false sharing、线程绑定和内存带宽曲线，但不能验证远端内存放置。承认证据边界，比编造一个漂亮结论更专业。

\section{本章习题}

\begin{exercisebox}
习题一：把一个全局原子计数器、每线程计数器、每 NUMA 节点计数器画成三张数据共享图。分别标出哪些 cache line 被谁写，哪些阶段需要合并，读取全局值的成本在哪里。

习题二：设计一个 NUMA first-touch 实验。要求写清初始化线程、计算线程、绑定策略、输入规模、测量指标和预期现象。若当前设备没有 NUMA，写出哪些部分只能作为计划保留。

习题三：为线程池设计一个两级队列方案：worker 本地队列加全局溢出队列。说明高频路径、低频路径、锁保护对象、可能的窃取顺序和关闭语义。

习题四：阅读 Linux NUMA 策略或调度拓扑的一个窄入口。报告必须从一个用户态实验问题出发，不允许只复制函数名。
\end{exercisebox}

\begin{labbox}
实验：在支持 NUMA 的机器上，用 \cmd{numactl --hardware} 查看节点。写一个大数组初始化和求和程序，比较单线程初始化、多线程初始化、线程绑定和未绑定情况下的带宽。如果没有 NUMA 机器，也要用 \cmd{lscpu} 记录拓扑并解释限制。
\end{labbox}
